\subsection{Checklist: Institutional review board approval}
\label{app:checklist_irb}

\paragraph{Question.}
Does the paper describe potential risks incurred by study participants, whether
such risks were disclosed to the subjects, and whether Institutional Review
Board approvals or an equivalent approval/review were obtained?

\paragraph{Answer (Yes/No/NA).}
NA.

\paragraph{Justification.}
The reported work does not involve human subjects. The manuscript studies
sparse Koopman autoencoders on numerical dynamical-system trajectory data. The
main benchmark consists of procedurally generated two-dimensional multibasin
systems with known attractor geometry, and the additional forecasting benchmark
uses public chaotic-flow systems from the Dysts package. The observations are
simulated state vectors, generated trajectories, basin metadata used only for
benchmark evaluation, model checkpoints, and aggregate forecasting or
support-diagnostic metrics.

No part of the described experimental protocol recruits, observes, surveys,
intervenes on, or collects data from people. The experiments do not use
participant records, medical or clinical data, user logs, demographic
attributes, private text, images, audio, biometric measurements, or other
personally identifiable or sensitive human-subject data. A search of the current
manuscript and source/script/test files for human-subjects terms such as IRB,
institutional review, participant, patient, survey, consent, privacy, and
personally identifiable data did not identify any human-subjects component.
Institutional review board approval is therefore not applicable to the current
paper.

\paragraph{Caveats.}
This answer applies only to the current manuscript and current experimental
code. Benchmark basin labels, attractor coordinates, and basin counts are labels
for simulated dynamical systems, not labels for people or protected groups. The
external Dysts dependency should still be disclosed, cited, and checked for
licensing requirements, but its use does not by itself create a human-subjects
review requirement.

If a later revision adds a user study, human-generated data, medical or
clinical measurements, user logs, private records, demographic attributes,
biometric data, surveys, interviews, safety-critical deployment with observed
people, or any dataset whose subjects are people, this checklist answer should
be changed and the authors should obtain and report the appropriate human-
subjects review, consent, privacy, and data-governance approvals before
submission.
